%% Cover letter for the Knowledge-Based Systems submission of
%% "HMRS: Content-informed hierarchical document management for scoped,
%% source-aware retrieval".
%% Scope follows Elsevier's cover-letter guidance: under one page; aim and
%% main findings; fit with the journal's aims and scope; novelty and wider
%% implications. Funding, author declarations and reviewer suggestions are
%% entered in Editorial Manager and are not repeated here.
%% AUTHOR ACTION: confirm the all-author approval, originality and
%% no-concurrent-submission statements before upload.
%% A line that starts with "[" must follow "\\{}" (else LaTeX reads it as
%% the optional argument of \\).
\documentclass[11pt,a4paper]{article}
\usepackage[margin=1in]{geometry}
\usepackage[T1]{fontenc}
\usepackage[utf8]{inputenc}
\usepackage{newtxtext}
\usepackage{enumitem}
\usepackage[hidelinks]{hyperref}
\setlength{\parindent}{0pt}
\setlength{\parskip}{3pt}
\pagestyle{empty}

\begin{document}

\textbf{Liangbin Li, Ph.D.}\\
{}National Synchrotron Radiation Laboratory, University of Science and Technology of China\\
{}Hefei, Anhui 230029, China\\
Email: lbli@ustc.edu.cn

\today

The Editors-in-Chief\\
Knowledge-Based Systems

Dear Editors,

\textbf{Re: ``HMRS: Content-informed hierarchical document management for scoped, source-aware retrieval''}

On behalf of all authors, I submit this manuscript for consideration as a
regular paper in Knowledge-Based Systems. The work is original, has not
been published elsewhere and is not under consideration by another
journal; all authors have approved this submission.

In changing, multi-domain document collections, a passage that resembles a
query need not come from a source relevant to the task. HMRS, the Hierarchical
Management Retrieval System, uses the same hierarchy for document management
and retrieval. An ingest agent uses recovered text to inform placement and
description. A query agent can consult the collection catalogue and document
listings before reading selected source passages, while lexical and vector
search remain available for candidate retrieval. Post-operation probes check
index visibility; an optional scorer admits passages but is not assumed to be
a suitable document-ranking score.

The paper's main findings are as follows.
\begin{itemize}[leftmargin=1.5em,itemsep=1pt,topsep=2pt,parsep=0pt]
\item A rename-and-move probe found 18 vector entries where nine were
expected; forced reindexing restored the expected count, but the probe did
not isolate the underlying fault.
\item In scripted tests on 72 questions, the reference route selected the
labelled collection on all 72, while the target string appeared in only 37
evidence packs. A multi-component configuration reached 52 target-string
hits and retained a higher count in two reruns on the same questions; these
are repeated observations on one cohort and do not isolate component effects.
\item In the fixed-pool replay, fused ordering kept gold documents in the top
five more often than verifier-score ordering. A separate native-deployment
comparison found no statistically significant paired difference after
correction; search scope and source access differed across deployments.
\end{itemize}

The manuscript presents a system-integration and observability contribution:
it links managed document state to agent-visible scope navigation and
post-mutation index checks, then distinguishes failures in placement, index
visibility, collection choice, text coverage and answer evaluation. It does
not claim a new ranking algorithm or a causal benefit over a tool-matched flat
retriever. The 72-question containment measure is a text-coverage proxy, not
claim-level support or answer accuracy.

Code and data are available at \url{https://github.com/kingdol666/rag-knowledge}.

The submission includes the \LaTeX{} source, highlights and supplementary
material. Thank you for considering our manuscript.

Yours sincerely,

\medskip
\textbf{Liangbin Li}, on behalf of all authors\\
{}National Synchrotron Radiation Laboratory, University of Science and Technology of China\\
{}Email: lbli@ustc.edu.cn

\end{document}
